\ProvidesFile{Locality.tex}[Полиморфизм и локальность]


\subsection{Полиморфизм и локальность}
\label{section::Locality}

В этом разделе я хочу поговорить про стоимость внедрения полиморфного поведения для классов.
Один из важных показателей -- на сколько локализованы изменения в коде.
Это в частности влияет и на количество кода, и на количество рутинных изменений в коде, и на количество файлов для перекомпиляции и линковки.
Я хочу разобрать конкретный пример для того, чтобы проследить что будет при классическом подходе и понять, что бы мы хотели в <<идеальном>> мире.
Начнем с простого сеттинга: у нас есть класс \verb"A" и функция \verb"func", которая принимает ссылку на объект класса \verb"A" и выполняет какие-то действия, и в \verb"main" мы зовем функцию на одном объекте.
Чтобы приблизить ситуацию к реальной, я разобью все на несколько файлов, как оно должно было бы быть в реальном проекте.
Это больше нужно, чтобы понять как добавление полиморфизма будет влиять на зависимости между файлами и на процесс компиляции.
{
\ProvidesFile{LocalityCode.tex}[Код для раздела Полиморфизм и локальность]


\newcommand{\codewidth}{5cm}

\newsavebox\funch
\begin{lrbox}{\funch}%
\begin{minipage}[\baselineskip]{\codewidth}
\texttt{function.h}
\begin{cppcode}[numbers = none]
#pragma once
#include "A.h"


void func(A&);

\end{cppcode}
\end{minipage}%
\end{lrbox}

\newsavebox\funccpp
\begin{lrbox}{\funccpp}%
\begin{minipage}[\baselineskip]{\codewidth}
\texttt{function.cpp}
\begin{cppcode}[numbers = none]
#include "function.h"

#include <iostream>

void func(A& a) {
  int x = a.f();
  int y = a.g(x);
  std::cout << y;
}
\end{cppcode}
\end{minipage}%
\end{lrbox}

\newsavebox\funcha
\begin{lrbox}{\funcha}%
\begin{minipage}[\baselineskip]{\codewidth}
\texttt{function.h}
\begin{cppcode}[numbers = none]
#pragma once
#include "???.h"


void func(???&);

\end{cppcode}
\end{minipage}%
\end{lrbox}

\newsavebox\funccppa
\begin{lrbox}{\funccppa}%
\begin{minipage}[\baselineskip]{\codewidth}
\texttt{function.cpp}
\begin{cppcode}[numbers = none]
#include "function.h"

#include <iostream>


void func(???& a) {
  int x = a.f();
  int y = a.g(x);
  std::cout << y;
}

\end{cppcode}
\end{minipage}%
\end{lrbox}

\newsavebox\funchb
\begin{lrbox}{\funchb}%
\begin{minipage}[\baselineskip]{\codewidth}
\texttt{function.h}
\begin{cppcode}[numbers = none]
#pragma once
#include "I.h"


void func(I&);

\end{cppcode}
\end{minipage}%
\end{lrbox}

\newsavebox\funccppb
\begin{lrbox}{\funccppb}%
\begin{minipage}[\baselineskip]{\codewidth}
\texttt{function.cpp}
\begin{cppcode}[numbers = none]
#include "function.h"

#include <iostream>

void func(I& a) {
  int x = a.f();
  int y = a.g(x);
  std::cout << y;
}
\end{cppcode}
\end{minipage}%
\end{lrbox}

\newsavebox\funccppc
\begin{lrbox}{\funccppc}%
\begin{minipage}[\baselineskip]{\codewidth}
\texttt{function.cpp}
\begin{cppcode}[numbers = none]
#include "function.h"

#include <iostream>


void func(I& a) {
  int x = a.f();
  int y = a.g(x);
  std::cout << y;
}

\end{cppcode}
\end{minipage}%
\end{lrbox}

\newsavebox\main
\begin{lrbox}{\main}%
\begin{minipage}[\baselineskip]{\codewidth}
\texttt{main.cpp}
\begin{cppcode}[numbers = none]
#include "function.h"

int main() {
  A a;
  func(a);
  A b;
  func(b);
  return 0;
}
\end{cppcode}
\end{minipage}%
\end{lrbox}

\newsavebox\maina
\begin{lrbox}{\maina}%
\begin{minipage}[\baselineskip]{\codewidth}
\texttt{main.cpp}
\begin{cppcode}[numbers = none]
#include "function.h"
#include "A.h"
#include "B.h"

int main() {
  A a;
  func(a);
  B b;
  func(b);
  return 0;
}
\end{cppcode}
\end{minipage}%
\end{lrbox}

\newsavebox\mainb
\begin{lrbox}{\mainb}%
\begin{minipage}[\baselineskip]{\codewidth}
\texttt{main.cpp}
\begin{cppcode}[numbers = none]
#include "function.h"
#include "A.h"
#include "BI.h"

int main() {
  A a;
  func(a);
  B b;
  BI bi = BI{&b};
  func(bi);
}
\end{cppcode}
\end{minipage}%
\end{lrbox}

\newsavebox\mainc
\begin{lrbox}{\mainc}%
\begin{minipage}[\baselineskip]{\codewidth}
\texttt{main.cpp}
\begin{cppcode}[numbers = none]
#include "function.h"

int main() {
  A a;
  func(a);
  A b;
  func(b);
  return 0;
}
\end{cppcode}
\end{minipage}%
\end{lrbox}

\newsavebox\classAh
\begin{lrbox}{\classAh}%
\begin{minipage}[\baselineskip]{\codewidth}
\texttt{A.h}
\begin{cppcode}[numbers = none]
#pragma once

struct A {
  int f() const;
  int g(int) const;
};
\end{cppcode}
\end{minipage}%
\end{lrbox}

\newsavebox\classAcpp
\begin{lrbox}{\classAcpp}%
\begin{minipage}[\baselineskip]{\codewidth}
\texttt{A.cpp}
\begin{cppcode}[numbers = none]
#include "A.h"

int A::f() const {
  return /*something*/;
}

int A::g(int x) const {
  return /*something of x*/;
}
\end{cppcode}
\end{minipage}%
\end{lrbox}

\newsavebox\classBh
\begin{lrbox}{\classBh}%
\begin{minipage}[\baselineskip]{\codewidth}
\texttt{B.h}
\begin{cppcode}[numbers = none]
#pragma once

struct B {
  int f() const;
  int g(int) const;
};
\end{cppcode}
\end{minipage}%
\end{lrbox}

\newsavebox\classBcpp
\begin{lrbox}{\classBcpp}%
\begin{minipage}[\baselineskip]{\codewidth}
\texttt{B.cpp}
\begin{cppcode}[numbers = none]
#include "B.h"

int B::f() const {
  return /*something*/;
}

int B::g(int x) const {
  return /*something of x*/;
}
\end{cppcode}
\end{minipage}%
\end{lrbox}

\newsavebox\classIh
\begin{lrbox}{\classIh}%
\begin{minipage}[\baselineskip]{\codewidth}
\texttt{I.h}
\begin{cppcode}[numbers = none]
#pragma once

struct I {
  int f() const;
  int g(int) const;
};
\end{cppcode}
\end{minipage}%
\end{lrbox}

\newsavebox\classAIh
\begin{lrbox}{\classAIh}%
\begin{minipage}[\baselineskip]{\codewidth}
\texttt{A.h}
\begin{cppcode}[numbers = none]
#pragma once

struct A : I {
  int f() const override;
  int g(int) const override;
};
\end{cppcode}
\end{minipage}%
\end{lrbox}

\newsavebox\classAIcpp
\begin{lrbox}{\classAIcpp}%
\begin{minipage}[\baselineskip]{\codewidth}
\texttt{A.cpp}
\begin{cppcode}[numbers = none]
#include "A.h"

int A::f() const {
  return /*something*/;
}

int A::g(int x) const {
  return /*something of x*/;
}
\end{cppcode}
\end{minipage}%
\end{lrbox}

\newsavebox\classBIh
\begin{lrbox}{\classBIh}%
\begin{minipage}[\baselineskip]{\codewidth}
\texttt{B.h}
\begin{cppcode}[numbers = none]
#pragma once

struct B : I {
  int f() const override;
  int g(int) const override;
};
\end{cppcode}
\end{minipage}%
\end{lrbox}

\newsavebox\classBIcpp
\begin{lrbox}{\classBIcpp}%
\begin{minipage}[\baselineskip]{\codewidth}
\texttt{B.cpp}
\begin{cppcode}[numbers = none]
#include "B.h"

int B::f() const {
  return /*something*/;
}

int B::g(int x) const {
  return /*something of x*/;
}
\end{cppcode}
\end{minipage}%
\end{lrbox}

\newsavebox\classIVh
\begin{lrbox}{\classIVh}%
\begin{minipage}[\baselineskip]{\codewidth+0.5cm}
\texttt{I.h}
\begin{cppcode}[numbers = none]
#pragma once

struct I {
  virtual int f() const = 0;
  virtual int g(int) const = 0;
};
\end{cppcode}
\end{minipage}%
\end{lrbox}

\newsavebox\classBIAh
\begin{lrbox}{\classBIAh}%
\begin{minipage}[\baselineskip]{\codewidth}
\texttt{BI.h}
\begin{cppcode}[numbers = none]
#pragma once
#include "I.h"
#include "B.h"

struct BI : I {
  int f() const override;
  int g(int) const override;
  B* b;
};
\end{cppcode}
\end{minipage}%
\end{lrbox}

\newsavebox\classBIAcpp
\begin{lrbox}{\classBIAcpp}%
\begin{minipage}[\baselineskip]{\codewidth}
\texttt{BI.cpp}
\begin{cppcode}[numbers = none]
#include "BI.h"

int B::f() const {
  return b->f();
}

int B::g(int x) const {
  return b->g(x);
}
\end{cppcode}
\end{minipage}%
\end{lrbox}


\paragraph{Исходный сеттинг}

Теперь у нас в нашем воображаемом проекте есть следующие файлы:
\[
\xymatrix{
  *+[F:<3pt>]{\usebox{\classAh}}&*+[F:<3pt>]{\usebox{\funch}}\ar[l]&{}\\
  *+[F:<3pt>]{\usebox{\classAcpp}}\ar[u]&*+[F:<3pt>]{\usebox{\funccpp}}\ar[u]&*+[F:<3pt>]{\usebox{\main}}\ar[ul]\\
}
\]

\paragraph{Добавление класса}

Теперь предположим, что что мы хотим добавить еще один класс \verb"B" с таким же интерфейсом и использовать его в функции \verb"func".
Теперь нас есть два почти идентичных класса:
\[
\xymatrix{
  *+[F:<3pt>]{\usebox{\classAh}}&*+[F:<3pt>]{\usebox{\classBh}}\\
  *+[F:<3pt>]{\usebox{\classAcpp}}\ar[u]&*+[F:<3pt>]{\usebox{\classBcpp}}\ar[u]
}
\]
И теперь использующий код будет выглядеть так
\[
\xymatrix{
  *+[F:<3pt>]{\usebox{\funcha}}&{}\\
  *+[F:<3pt>]{\usebox{\funccppa}}\ar[u]&*+[F:<3pt>]{\usebox{\maina}}\ar[ul]\\
}
\]
И вот теперь надо изменить сигнатуру функции, чтобы изменить код.
Добавление класса \verb"B" -- это локальное действие, потому что мы никак не влияем ни на какие файлы кроме \verb"B.h" и \verb"B.cpp".
Изменение в функции \verb"main" все еще локальные, мы хотели добавить в нее поддержку \verb"B", вот мы и изменили этот код.
Теперь нужно лишь понять, как изменить \verb"func", чтобы в нее можно было подставлять и объекты класса \verb"A" и объекты класса \verb"B".
Для этого нужно завести новый класс \verb"I" являющийся интерфейсом для \verb"A" и \verb"B", через который \verb"func" будет взаимодействовать с классами.
То есть в идеале мы бы хотели что-то такое%
\footnote{На всякий случай скажу, что код ниже не рабочий.
Это пока рассуждения, что мы хотели бы сделать.}
\[
\xymatrix{
  *+[F:<3pt>]{\usebox{\classIh}}&*+[F:<3pt>]{\usebox{\funchb}}\ar[l]&{}\\
  {}&*+[F:<3pt>]{\usebox{\funccppb}}\ar[u]&*+[F:<3pt>]{\usebox{\mainc}}\ar[ul]\\
}
\]
И вот теперь остается вопрос, а как прикрепить этот самый класс \verb"I" с интерфейсом к классам \verb"A" и \verb"B".
И вот тут есть два пути:
\begin{enumerate}
\item Классический подход через наследование с виртуальными функциями.

\item Подход через стирающие типы и ссылки (указатели).
\end{enumerate}
Я хочу сейчас проследить на сколько классический подход не является локальным.
На сколько много изменений он требует в коде, в структуре компиляции проекта, и как много потребуется лишней рутинной работы.

\paragraph{Классическое решение}

Добавление интерфейса к классу происходит путем наследования от интерфейса.
При этом методы интерфейса надо сделать виртуальными.
А все значит, что у каждого класса, который мы хотим использовать полиморфно надо внести эти изменения в файлы класса.
\[
\xymatrix@C=10pt{
  {}&*+[F:<3pt>]{\usebox{\classIVh}}&{}\\
  *+[F:<3pt>]{\usebox{\classAIh}}\ar[ur]&{}&*+[F:<3pt>]{\usebox{\classBIh}}\ar[ul]\\
  *+[F:<3pt>]{\usebox{\classAIcpp}}\ar[u]&{}&*+[F:<3pt>]{\usebox{\classBIcpp}}\ar[u]\\
}
\]
Отметим, что на уровне используемых классов данное решение интрузивное.
То есть оно вмешивается во все используемые классы \verb"A" и \verb"B".
Но что если вам нельзя изменить код класса \verb"B"?
Что если он является каким-то сторонним классом из другой библиотеки?
Тогда вам придется писать адаптер следующего вида.
\[
\xymatrix{
  *+[F:<3pt>]{\usebox{\classIVh}}&*+[F:<3pt>]{\usebox{\classBIAh}}\ar[l]\ar[r]&*+[F:<3pt>]{\usebox{\classBh}}\\
  {}&*+[F:<3pt>]{\usebox{\classBIAcpp}}\ar[u]&*+[F:<3pt>]{\usebox{\classBcpp}}\ar[u]\\
}
\]
Давайте поясню, здесь класс \verb"BI" унаследован от интерфейса, а значит может быть использован как аргумент функции \verb"func".
А объект класса \verb"BI" хранит в себе указатель на объект класса \verb"B" и просто пробрасывает методы наружу.
То есть везде, где надо использовать класс \verb"B" вам придется использовать обертку \verb"BI".
И к сожалению в случае функции \verb"func" нельзя будет использовать временный объект класса \verb"BI".

Если у нас не было сторонних классов, то теперь использующий код будет выглядеть следующим образом%
\footnote{Вот теперь этот код будет рабочим.
Хотя сделать виртуальным деструктор тут бы все же не мешало.}
\[
\xymatrix{
  *+[F:<3pt>]{\usebox{\classIVh}}&*+[F:<3pt>]{\usebox{\funchb}}\ar[l]&{}\\
  {}&*+[F:<3pt>]{\usebox{\funccppc}}\ar[u]&*+[F:<3pt>]{\usebox{\maina}}\ar[ul]\\
}
\]
Если же класс \verb"B" принадлежал сторонней библиотеке, то нам придется все же изменить код \verb"main" следующим образом:
\[
\xymatrix{
  *+[F:<3pt>]{\usebox{\classIVh}}&*+[F:<3pt>]{\usebox{\funchb}}\ar[l]&*+[F:<3pt>]{\usebox{\classBIAh}}\\
  {}&*+[F:<3pt>]{\usebox{\funccppc}}\ar[u]&*+[F:<3pt>]{\usebox{\mainb}}\ar[ul]\ar[u]\\
}
\]
То есть код на стороне пользователя просто так не останется неизменным.
Ситуация становится сильно хуже, если мы хотим не просто ссылаться на полиморфный объект, как это сделано в примере выше, а если мы хотим владеть им.
Например передать полиморфный объект внутрь и потом вернуть его наружу с помощью мув операций.
Или принять его копию по значению для пользования.
Или же завладеть им, для размещения в контейнере.

Давайте сейчас сравним диаграммы зависимостей в файлах до и после полиморфизма.
До полиморфизма ситуация такая
\[
\xymatrix{
  *+[F:<3pt>]{\texttt{A.h}}&*+[F:<3pt>]{\texttt{function.h}}\ar[l]&{}\\
  *+[F:<3pt>]{\texttt{A.cpp}}\ar[u]&*+[F:<3pt>]{\texttt{function.cpp}}\ar[u]&*+[F:<3pt>]{\texttt{main.cpp}}\ar[ul]\\
}
\]
И после добавления полиморфизма (считаем, что \verb"B" внешний класс, в который нельзя вмешиваться).
\[
\colorlet{lightred}{red!30}
\xymatrix{
  {}&{}&{}&*+[F:<3pt>]{\texttt{I.h}}&{}\\
  *+[F:<3pt>]{\texttt{B.h}}&
  *+[F:red:<3pt>]{\texttt{BI.h}}\ar@[red][l]\ar@[red][urr]&
  *+[F:<3pt>][F*:lightred:<3pt>]{\texttt{A.h}}\ar@[red][ur]&
  {}&
  *+[F:<3pt>][F*:lightred:<3pt>]{\texttt{function.h}}\ar[ul]\\
  *+[F:<3pt>]{\texttt{B.cpp}}\ar[u]&
  *+[F:red:<3pt>]{\texttt{BI.cpp}}\ar@[red][u]&
  *+[F:<3pt>]{\texttt{A.cpp}}\ar[u]&
  *+[F:<3pt>][F*:lightred:<3pt>]{\texttt{main.cpp}}\ar[ur]\ar[ul]\ar@[red][ull]&
  *+[F:<3pt>]{\texttt{function.cpp}}\ar[u]\\
}
\]
Понятно, что часть файлов и зависимостей обязана измениться, просто потому что мы добавили еще один класс \verb"B".
Так же у нас появился общий интерфейс и это тоже должно отразится на коде.
Однако, у нас появилось много лишних и не нужных зависимостей.
Я отметил такие зависимости красным.
\begin{enumerate}
\item Появился файл для интерфейса \verb"I.h".
Это неизбежно.

\item Файл \verb"A.h" теперь зависит от \verb"I.h".

\item Появился лишний класс адаптер \verb"BI.h/cpp".

\item Файл \verb"function.h" теперь зависит от интерфейса \verb"I.h", а не от конкретного файла класса.
Это неизбежно.

\item \verb"main.cpp" зависит не от файла \verb"B", а от файле для его адаптера \verb"BI.h"
\end{enumerate}
Так же красной заливкой я отметил существующие файлы, которые требуют изменения.
\begin{enumerate}
\item Да код \verb"main.cpp" мы ожидаем, что будет другой, потому что мы туда хотим добавить новый класс.
Однако, мы вместо этого добавляем туда зависимость от целого адаптера, а не от нового класса.

\item Код для \verb"funcion.h" почти не меняется и добавление зависимости от интерфейса \verb"I.h" это ожидаемое решение.

\item Код для \verb"A.h" мы вынуждены поменять, так как классическое решение интрузивное и вмешивается в код класса.
\end{enumerate}
Что происходит с единицами компиляции.
\begin{enumerate}
\item Файл \verb"function.cpp" надо перекомпилировать и от этого не избавиться, так как мы меняем сигнатуру функции.

\item Файл \verb"main.cpp" надо перекомпилировать и от этого не избавиться, потому что мы хотим писать туда новый код.

\item Файл \verb"A.cpp" надо перекомпилировать.
И это лишняя работа, так как у нас нет необходимости вносить изменения в сам класс.
Тем не менее, мы наследованием изменили layout самого класса и добавили виртуальных функций.

\item Файл \verb"B.cpp" считается внешним и его менять нельзя.
В данном случае эта единица трансляции не меняется.

\item Появилась новая единица трансляции для адаптера \verb"BI.cpp".
Это вынужденный костыль, никак по-другому сторонний класс не прицепишь.
\end{enumerate}
По хорошему, мы бы хотели иметь следующие изменения:
\[
\colorlet{lightred}{red!30}
\xymatrix{
  *+[F:<3pt>]{\texttt{B.h}}&
  *+[F:<3pt>]{\texttt{A.h}}&
  *+[F:red:<3pt>]{\texttt{I.h}}&
  *+[F:<3pt>][F*:lightred:<3pt>]{\texttt{fucnction.h}}\ar@[red][l]\\
  *+[F:<3pt>]{\texttt{B.cpp}}\ar[u]&
  *+[F:<3pt>]{\texttt{A.cpp}}\ar[u]&
  *+[F:<3pt>][F*:lightred:<3pt>]{\texttt{main.cpp}}\ar@[red][ul]\ar@[red][ull]\ar[ur]&
  *+[F:<3pt>]{\texttt{function.cpp}}\ar[u]\\
}
\]
Давайте отметим ключевые вещи (все описанные изменения неизбежны):
\begin{enumerate}
\item Мы добавляем только один новый файл с описанием интерфейса \verb"I.h".

\item Меняем заголовок \verb"function.h" и добавляем зависимость от интерфейса \verb"I.h".

\item Добавляем напрямую зависимости в \verb"main.cpp" от файлов классов \verb"A.h" и \verb"B.h".

\item Меняем бизнес логику в коде \verb"main.cpp".

\item Мы не добавляем лишних единиц трансляции.

\item Перекомпилируются только те единицы трансляции, которые вынуждены будут перекомпилироваться, а именно \verb"function.cpp" и \verb"main.cpp".
\end{enumerate}
Следующий раздел как раз и будет посвящен решению, которое удовлетворяет указанной выше схеме.
Это решение основано на стирающих типах, ссылках и указателях.

\paragraph{Владение и классическое решение}

Давайте посмотрим, что будет 
% TO DO
% отправить сообщение в MDS и владеть неизвестными данными

}
